%%%%%%%%%%%%%%%%%%%%%%%%%%%%%%%%%%%%%%%%%%%%%%%%%%%%%%%%%%%%%%%%%%%%%%
% 5. CONCLUSÕES
%%%%%%%%%%%%%%%%%%%%%%%%%%%%%%%%%%%%%%%%%%%%%%%%%%%%%%%%%%%%%%%%%%%%%%
\section{CONCLUSÕES}

Este trabalho apresentou o desenvolvimento e a validação preliminar do SRAB para PocketQube 1P. Destacam-se cinco contribuições principais: (1) o modelo dinâmico validou a relação não linear do torque aerodinâmico, evidenciada pelo aumento de 81,6\% na conicidade de impacto (de 13,46$^\circ$ para 24,44$^\circ$) na transição Asa1--Asa2; (2) a instrumentação embarcada foi funcionalmente validada em bancada para aquisição a 20~Hz; (3) a simulação da configuração Asa3 (4 asas, 1000~m) resultou em velocidade de impacto de 12,90~m/s, conicidade de 17,22$^\circ$ e \textit{spin} de 372~RPM, dissipando 99,15\% da energia potencial; (4) a velocidade terminal mostrou-se independente de variações de $\pm20\%$ no $C_{D0}$, confirmando a dominância do LEV e atestando a robustez para impressão 3D; e (5) ante um paraquedas de 200~mm, o SRAB reduziu o tempo de descida em 34\%, mitigando a deriva por vento e eliminando mecanismos de acionamento passíveis de falha.

A análise Monte Carlo (200 iterações) atestou a alta estabilidade do sistema frente a incertezas paramétricas (massa, $\beta$, $C_{D0}$, $f_{factor}$), estimando uma velocidade de impacto de 12,89~$\pm$~0,20~m/s (CV~$<$~2\%). Embora a configuração atual opere de forma mais lenta que a exigida pela janela regulatória LASC (20--45~m/s), demandando ajustes geométricos simples, como alteração no número de asas ou no ângulo de passo $\beta$, a baixa energia cinética de impacto reforça a segurança do conceito.

Conclui-se que o SRAB demonstra plena viabilidade como sistema de recuperação passivo para nanossatélites, destacando-se pela extrema confiabilidade estrutural. Os trabalhos futuros buscarão a transição da prova de conceito para a certificação operacional.

%%%%%%%%%%%%%%%%%%%%%%%%%%%%%%%%%%%%%%%%%%%%%%%%%%%%%%%%%%%%%%%%%%%%%%
% AGRADECIMENTOS
%%%%%%%%%%%%%%%%%%%%%%%%%%%%%%%%%%%%%%%%%%%%%%%%%%%%%%%%%%%%%%%%%%%%%%
\subsection*{\textit{Agradecimentos}}

Os autores agradecem à equipe Serra Rocketry (Missão Helike \#213), à Alba Orbital pela doação da plataforma PocketQube 1P e aos laboratórios do IPRJ/UERJ pelo suporte institucional.